\section{Introduction}

Teacher bias may be central to inequalities in educational attainment and later life outcomes. To estimate the extent of such bias, researchers typically compare student assessments in which teachers exercise discretion, such as teacher-assigned grades or track recommendations, to more objective indicators of student ability, such as standardized test scores \citep{alesina2024revealing,falk2023mentoring}. However, the presence of systematic gaps in subjective teacher evaluations between groups, conditional on objective ability measures, is not necessarily evidence of teacher bias. This approach faces two main concerns: omitted variables and measurement error in the ability measure. Studies often address the first concern by controlling for additional factors \citep{burgess2013test}. In contrast, the second concern is typically overlooked or addressed under strict assumptions that may not hold in practice, a shortcoming shared by the broader literature on discrimination.

This paper focuses on the second concern: measurement error in the objective measure of ability. In classical testing theory, test scores are typically based on the number of correct answers and provide an unbiased but noisy estimate of ability. The resulting ``classical'' measurement error in the test score is considered purely random and hence uncorrelated with ability. In modern item response theory (IRT), statistical models produce test scores that may be biased estimates of ability but are unbiased as predictors. The resulting ``non-classical'' measurement error may correlate with ability. For instance, in case the test score is estimated with a posterior mean score, students with a low and high signal are drawn towards the population mean, such that the error correlates negatively with ability, yet is uncorrelated with the test score. 

To examine the implications, consider a regression in which a teacher’s subjective evaluation of a student is the dependent variable, and the independent variables include a group characteristic, such as family socioeconomic status (SES), and a test score as a measure of student ability. In case of classical measurement error, the test score coefficient is biased towards zero (the well-known attenuation bias). Moreover, if student ability is associated with SES and matters for teacher evaluations, this bias also contaminates the coefficient of the conditional SES gap. In particular, when ability is positively associated with both SES and teacher evaluations, the conditional SES gap is overestimated. In case of non-classical measurement error, the attenuation bias in the test score coefficient may be partially or fully eliminated, though this leads to a biased coefficient of the association between ability and SES. This bias then also contaminates the estimate for the conditional SES gap.

We first show how both classical and certain forms of non-classical measurement error, particularly relevant to modern test scores, generate identical contamination bias in the coefficient measuring the conditional gap between groups. We then discuss three empirical strategies to address measurement error. First, we consider an instrumental variable (IV) approach, a well-established method in econometrics to correct for measurement error \citep{hausman2001mismeasured}. Second, we discuss a method in the spirit of errors-in-variables (EIV) models, an approach that is common in psychometrics \citep{fuller2009measurement}. The additional parameter required for the EIV strategy, compared to OLS, is the reliability ratio of the test that is used to measure ability. We argue that, under weaker assumptions than the IV strategy, the IV first stage can be used to identify this ratio. Third, we combine the EIV strategy with a novel approach to identify the reliability ratio under yet weaker assumptions. 

We apply these three strategies to estimate conditional gaps in teacher track recommendations in the Netherlands, focusing on differences by family SES. The Dutch tracking regime provides a compelling setting for this analysis. While nearly all OECD countries have some form of educational tracking, the Netherlands assigns students to secondary-school tracks at a relatively young age of twelve. Assignment is based on binding track recommendations from the primary school teacher \citep{wvo2014}. These recommendations may be highly consequential, as roughly 70\% of students does not experience track mobility four years into secondary education \citep{deree2023quality}. Track placement determines the level of the curriculum and peer group abilities, and may affect student achievement \citep{matthewes2021better}. As only the highest tracks provide access to university, track placement may also have long-term consequences for attainment and income \citep{dustmann2017long,borghans2019long}. Moreover, negatively biased recommendations signal low teacher expectations, which can in turn harm student outcomes \citep{carlana2019implicit,papageorge2020teacher}. 

Given the consequential nature of these recommendations, potential teacher bias has been a major policy concern in the Netherlands. This concern gained momentum in 2016, when a report by the Dutch \citet{inspectorate2016} demonstrated that students from low SES families receive significantly lower track recommendations, conditional on scores from a standardized end-of-primary school test. This finding fueled an ongoing public debate on teachers as gatekeepers of opportunities and contributed to a national policy reform, effective from 2023-24. The reform reduces the role of teachers in track placement, and increases the role of the end-of-primary school test. Similar findings in academic studies reinforce concerns about unequal opportunities in education, both in the Netherlands \citep{zumbuehl2025can} and elsewhere \citep{carlana2022implicit,falk2023mentoring}.

Using new administrative data, we first replicate previous findings from the Netherlands, and find that children with lower family SES receive systematically lower track recommendations, conditional on standardized test scores. We discuss that our test scores are derived from an IRT model, with ability parameters estimated using weighted maximum likelihood \citep{sanders1993psychometrie}. This method is known to produce largely unbiased estimates of student ability \citep{warm1989weighted}, and the measurement error is considered classical. We address concerns about measurement error using three strategies. Our novel strategy combines the EIV method with an unbiased estimate for the test's reliability ratio under weak assumptions, based on students’ standardized test scores across all primary school grades. Across the three approaches, we find that measurement error can explain between 35 and 43\% of the conditional SES gap. 

Our paper contributes to the broad literature on measurement error \citep{schennach2016recent}. In particular, we focus on contamination bias due to measurement error. As \citet{modalsli2022spillover} point out,
%
``\textit{[a]lthough the notion of bias in one coefficient arising from error in another regressor is a well-known econometric result, it is seldom addressed in practice with empirical studies.}''
%
\citet{modalsli2022spillover} examine the implications of classical measurement error in multigenerational income regressions and use an IV approach to mitigate the resulting bias. Similarly, \citet{gillen2019experimenting} use IV to show that the gender gap in competitiveness is well explained by risk attitudes and overconfidence once classical measurement error is addressed. Our study extends this literature by applying three strategies to address contamination bias and by incorporating both classical and non-classical measurement error. The latter is pervasive in education, where test scores are often derived from IRT models \citep{jacob2016measurement}. 
It is also common in the labor market, as many individuals misreport hours worked, leading to biased estimates of wage discrimination and inequality \citep{borjas2024}. More generally, it arises in survey data when respondents base their answers on incomplete information \citep{hyslop2001bias}

We also contribute to the literature on teacher bias by providing new evidence on the conditional SES gap in track recommendations. Teacher bias in track recommendations has been studied in Germany \citep{falk2023mentoring}, Italy \citep{carlana2022implicit}, and the Netherlands \citep{zumbuehl2025can}. Previous studies have also examined teacher bias in grading by gender \citep{cornwell2013noncognitive,lavy2018origins,terrier2020boys}, ethnicity \citep{burgess2013test}, and migration background \citep{alesina2024revealing}. Consistent with the observation above, contamination bias due to measurement error has been largely ignored in this literature. Three notable exceptions are \citet{botelho2015racial}, \cite{ferman2022assessing}, and \cite{zhu2024}, who study teacher bias in grading in Brazil and the US. All three studies find that the conditional gaps are substantially reduced or even reversed when addressing measurement error in test scores. However, as these studies assume classical measurement error and exclusively apply an IV strategy, it remains unclear whether the results hold under non-classical error and are robust to alternative strategies relying on weaker assumptions.

